\documentclass{beamer}
\mode<presentation>
\usepackage{amsmath,amssymb,mathtools}
\usepackage{textcomp}
\usepackage{gensymb}
\usepackage{adjustbox}
\usepackage{subcaption}
\usepackage{enumitem}
\usepackage{multicol}
\usepackage{listings}
\usepackage{url}
\usepackage{graphicx} % <-- needed for images
\def\UrlBreaks{\do\/\do-}

\usetheme{Boadilla}
\usecolortheme{lily}
\setbeamertemplate{footline}{
  \leavevmode%
  \hbox{%
  \begin{beamercolorbox}[wd=\paperwidth,ht=2ex,dp=1ex,right]{author in head/foot}%
    \insertframenumber{} / \inserttotalframenumber\hspace*{2ex}
  \end{beamercolorbox}}%
  \vskip0pt%
}
\setbeamertemplate{navigation symbols}{}

\lstset{
  frame=single,
  breaklines=true,
  columns=fullflexible,
  basicstyle=\ttfamily\tiny   % tiny font so code fits
}

\numberwithin{equation}{section}

% ---- your macros ----
\providecommand{\nCr}[2]{\,^{#1}C_{#2}}
\providecommand{\nPr}[2]{\,^{#1}P_{#2}}
\providecommand{\mbf}{\mathbf}
\providecommand{\pr}[1]{\ensuremath{\Pr\left(#1\right)}}
\providecommand{\qfunc}[1]{\ensuremath{Q\left(#1\right)}}
\providecommand{\sbrak}[1]{\ensuremath{{}\left[#1\right]}}
\providecommand{\lsbrak}[1]{\ensuremath{{}\left[#1\right.}}
\providecommand{\rsbrak}[1]{\ensuremath{\left.#1\right]}}
\providecommand{\brak}[1]{\ensuremath{\left(#1\right)}}
\providecommand{\lbrak}[1]{\ensuremath{\left(#1\right.}}
\providecommand{\rbrak}[1]{\ensuremath{\left.#1\right)}}
\providecommand{\cbrak}[1]{\ensuremath{\left\{#1\right\}}}
\providecommand{\lcbrak}[1]{\ensuremath{\left\{#1\right.}}
\providecommand{\rcbrak}[1]{\ensuremath{\left.#1\right\}}}
\theoremstyle{remark}
\newtheorem{rem}{Remark}
\newcommand{\sgn}{\mathop{\mathrm{sgn}}}
\providecommand{\abs}[1]{\left\vert#1\right\vert}
\providecommand{\res}[1]{\Res\displaylimits_{#1}}
\providecommand{\norm}[1]{\lVert#1\rVert}
\providecommand{\mtx}[1]{\mathbf{#1}}
\providecommand{\mean}[1]{E\left[ #1 \right]}
\providecommand{\fourier}{\overset{\mathcal{F}}{ \rightleftharpoons}}
\providecommand{\system}{\overset{\mathcal{H}}{ \longleftrightarrow}}
\providecommand{\dec}[2]{\ensuremath{\overset{#1}{\underset{#2}{\gtrless}}}}
\newcommand{\myvec}[1]{\ensuremath{\begin{pmatrix}#1\end{pmatrix}}}
\newcommand{\mydet}[1]{\ensuremath{\begin{vmatrix}#1\end{vmatrix}}}

\newenvironment{amatrix}[1]{%
  \left(\begin{array}{@{}*{#1}{c}|*{#1}{c}@{}}
}{%
  \end{array}\right)
}

\newcommand{\myaugvec}[2]{\ensuremath{\begin{amatrix}{#1}#2\end{amatrix}}}
\let\vec\mathbf
% ---------------------

\title{Eigen Values and Vectors}
\author{ee25btech11056 - Suraj.N}

\begin{document}

\begin{frame}
  \titlepage
\end{frame}

\begin{frame}{Problem Statement}

What can we conclude if a matrix $\vec{A}$ is symmetric?\\

\end{frame}

\begin{frame}{Proof}

\textbf{Proof :}\\

\textbf{Case 1 : $\vec{A}$ is real and symmetric}\\

Let $\vec{A} \in \mathbb{R}^{n \times n}$ be a real symmetric matrix, i.e.,
\begin{align}
\vec{A} = \vec{A}^\top
\end{align}

Suppose $\lambda$ is an eigenvalue of $\vec{A}$ with eigenvector $\vec{v} \ne \vec{0}$:
\begin{align}
\vec{A}\vec{v} = \lambda \vec{v}
\end{align}

Taking the conjugate transpose of $\vec{v}$ and multiplying on the left:
\begin{align}
\vec{v}^* \vec{A}\vec{v} = \lambda \vec{v}^* \vec{v}
\end{align}

Since $\vec{A}$ is real symmetric, $\vec{A} = \vec{A}^* = \vec{A}^\top$, therefore:
\begin{align}
\vec{v}^* \vec{A}\vec{v} = (\vec{A}\vec{v})^* \vec{v} = (\lambda \vec{v})^* \vec{v} = \bar{\lambda} \, \vec{v}^* \vec{v}
\end{align}

Equating both expressions for $\vec{v}^* \vec{A}\vec{v}$:
\begin{align}
\lambda \vec{v}^* \vec{v} = \bar{\lambda} \, \vec{v}^* \vec{v}\\
(\lambda- \bar{\lambda})\vec{v}^* \vec{v} = \vec{0}
\end{align}

\end{frame}

\begin{frame}{Proof}

Since $\vec{v} \ne \vec{0}$, $\vec{v}^*\vec{v} > 0$. Hence:
\begin{align}
\lambda = \bar{\lambda}
\end{align}

Therefore, every eigenvalue $\lambda$ of $\vec{A}$ is real.\\

\textbf{Eigenvectors corresponding to distinct eigenvalues are orthogonal}\\

Let

\begin{align}
\vec{A}\vec{u} = \alpha \vec{u}\\
\vec{A}\vec{w} = \beta \vec{w}
\end{align}

Where $\alpha \ne \beta$ are distinct eigenvalues.Multiplying the first equation on the left by $\vec{w}^\top$:

\begin{align}
\vec{w}^\top \vec{A} \vec{u} = \alpha \vec{w}^\top \vec{u}
\end{align}

\end{frame}

\begin{frame}{Proof}

Since $\vec{A}$ is symmetric, $\vec{A} = \vec{A}^\top$. Hence:
\begin{align}
\vec{w}^\top \vec{A}\vec{u} = (\vec{A}\vec{w})^\top \vec{u} = (\beta \vec{w})^\top \vec{u} = \beta \vec{w}^\top \vec{u}
\end{align}

Equating both forms:
\begin{align}
\alpha \vec{w}^\top \vec{u} = \beta \vec{w}^\top \vec{u}
\end{align}

Simplifying:
\begin{align}
(\alpha - \beta) \vec{w}^\top \vec{u} = 0
\end{align}

Since $\alpha \ne \beta$, it follows that:
\begin{align}
\vec{w}^\top \vec{u} = 0
\end{align}

Hence, eigenvectors corresponding to distinct eigenvalues are orthogonal.\\

\textbf{Proving that Eigen Vector is real}

Suppose $\lambda \in \mathbb{R}$ is an eigenvalue of $\vec{A}$ with a (possibly complex) eigenvector $\vec{v} \in \mathbb{C}^n$, satisfying
\begin{align}
\vec{A}\vec{v} = \lambda \vec{v}
\end{align}

\end{frame}

\begin{frame}{Proof}

\textbf{Considering the conjugate eigenvector}

Since $\vec{A}$ is real, taking the complex conjugate of both sides gives
\begin{align}
\vec{A}\bar{\vec{v}} = \bar{\lambda} \bar{\vec{v}}
\end{align}

But for a real symmetric matrix, $\lambda = \bar{\lambda}$ (as proven earlier). Hence,
\begin{align}
\vec{A}\bar{\vec{v}} = \lambda \bar{\vec{v}}
\end{align}

Thus, $\bar{\vec{v}}$ is also an eigenvector corresponding to the same eigenvalue $\lambda$.

\bigskip

\textbf{Constructing real vectors from $\vec{v}$ and $\bar{\vec{v}}$}

Define
\begin{align}
\vec{u} &= \vec{v} + \bar{\vec{v}}\\
\vec{w} &= i(\vec{v} - \bar{\vec{v}})
\end{align}

Both $\vec{u}$ and $\vec{w}$ are \textbf{real vectors}, because $\vec{u}$ is the sum of a vector and its conjugate, and $\vec{w}$ is $i$ times their difference.

\bigskip
\end{frame}

\begin{frame}{Proof}

\textbf{Applying $\vec{A}$ to $\vec{u}$ and $\vec{w}$}

Since $\vec{A}$ is real and linear,
\begin{align}
\vec{A}\vec{u} &= \vec{A}(\vec{v} + \bar{\vec{v}}) = \vec{A}\vec{v} + \vec{A}\bar{\vec{v}} = \lambda \vec{v} + \lambda \bar{\vec{v}} = \lambda (\vec{v} + \bar{\vec{v}}) = \lambda \vec{u}
\end{align}

Similarly,
\begin{align}
\vec{A}\vec{w} &= \vec{A}\big(i(\vec{v} - \bar{\vec{v}})\big) = i(\vec{A}\vec{v} - \vec{A}\bar{\vec{v}}) = i(\lambda \vec{v} - \lambda \bar{\vec{v}}) = \lambda i(\vec{v} - \bar{\vec{v}}) = \lambda \vec{w}
\end{align}

Hence we get 

\begin{align}
  \vec{A}\vec{u}=\lambda\vec{u}\\
  \vec{A}\vec{w}=\lambda\vec{w}
\end{align}
\bigskip

\textbf{Real eigenvectors}

If $\vec{u} \ne \vec{0}$, then $\vec{u}$ is a real eigenvector corresponding to eigenvalue $\lambda$.  
If $\vec{w} \ne \vec{0}$, then $\vec{w}$ is a real eigenvector corresponding to eigenvalue $\lambda$.  

\end{frame}

\begin{frame}{Proof}

\textbf{Conclusion :}
For any real symmetric matrix $\vec{A}$,

  1) All eigenvalues of $\vec{A}$ are \textbf{real}.\\
  2) Eigenvectors corresponding to distinct eigenvalues are \textbf{orthogonal}.\\
  3) For every eigenvalue $\lambda$ of a real symmetric matrix $\vec{A}$, there exists a \textbf{real eigenvector}.  
Hence, the eigenvectors of real symmetric matrices are real. 


\textbf{Case 2 : $\vec{A}$ is complex and symmetric}\\

Let $\vec{A} \in \mathbb{C}^{n \times n}$ be a complex symmetric matrix satisfying
\begin{align}
\vec{A} = \vec{A}^\top
\end{align}

Note that $\vec{A}$ is not necessarily Hermitian since in general $\vec{A} \ne \vec{A}^*$.

\begin{align}
\vec{A} \ne \vec{A}^*
\end{align}

\textbf{Eigenvalues of $\vec{A}$ may be complex}

Let $\lambda \in \mathbb{C}$ be an eigenvalue of $\vec{A}$ with eigenvector $\vec{v} \ne \vec{0}$ such that
\begin{align}
\vec{A}\vec{v} = \lambda \vec{v}
\end{align}

\end{frame}

\begin{frame}{Proof}

Define scalar z as :

\begin{align}
z = \vec{v}^* \vec{A} \vec{v}  
\end{align}

Taking complex conjugate transpose on both sides ,

\begin{align}
z^* = \vec{v}^* \vec{A}^* \vec{v}   
\end{align}

If $\vec{A} = \vec{A}^*$ , then 

\begin{align}
  z = z^* \\
  z = \bar{z} \text{   as z is a scalar}
\end{align}

Therefore z is a real scalar value \\

For Hermitian matrices ($\vec{A} = \vec{A}^*$), the quantity $\vec{v}^* \vec{A} \vec{v}$ is always real.  
However, for complex symmetric matrices, $\vec{A} \ne \vec{A}^*$, so

\end{frame}

\begin{frame}{Proof}

\begin{align}
\vec{v}^* \vec{A} \vec{v} \text{ may not be real.}
\end{align}

The Rayleigh quotient is
\begin{align}
\frac{\vec{v}^* \vec{A} \vec{v}}{\vec{v}^* \vec{v}} = \lambda
\end{align}

Since $\vec{v}^* \vec{A} \vec{v}$ may be complex, $\lambda$ may also be complex.

\textbf{Eigenvectors corresponding to distinct eigenvalues need not be orthogonal}

For complex symmetric matrices, we have only $\vec{A} = \vec{A}^\top$, so the Hermitian symmetry property does not hold.

Thus, for two eigenpairs $(\lambda_1, \vec{v}_1)$ and $(\lambda_2, \vec{v}_2)$ with $\lambda_1 \ne \lambda_2$, we \textbf{cannot guarantee}
\begin{align}
\vec{v}_1^* \vec{v}_2 = 0
\end{align}

\end{frame}

\begin{frame}{Proof}

Let

\begin{align}
\vec{A}\vec{v_1} = \alpha \vec{v_1}\\
\vec{A}\vec{v_2} = \beta \vec{v_2}
\end{align}

Where $\alpha \ne \beta$ are distinct eigenvalues.Multiplying the first equation on the left by $\vec{v_2}^*$:
\begin{align}
\vec{v_2}^* \vec{A} \vec{v_1} = \alpha \vec{v_2}^* \vec{v_1}
\end{align}

As $\vec{A} \neq \vec{A}^*$ and also $\vec{A} \neq \bar{\vec{A}}$ we cannot proceed like what we had done earlier for \textbf{case 1}.

Hence, eigenvectors corresponding to distinct eigenvalues \textbf{need not be orthogonal and are not guaranteed to be real}.

\textbf{Conclusion :}  
For a complex symmetric matrix $\vec{A}$:

  1) The eigenvalues \textbf{can be complex}.\\
  2) The eigenvectors corresponding to distinct eigenvalues are \textbf{not necessarily orthogonal and need not be real}.

 
But if $\vec{A} = \vec{A}^*$ then we get the same result as \textbf{case 1}

\end{frame}

\begin{frame}{Examples}

\textbf{Examples to Support the Proof} 

\textbf{Example 1 : Real and Symmetric Matrix }

(\textbf{a}) Let
\begin{align}
\vec{A} = \myvec{2 & -1 & 0\\ -1 & 2 & -1\\ 0 & -1 & 2}
\end{align}

Clearly, $\vec{A} = \vec{A}^\top$, hence $\vec{A}$ is real symmetric.

\textbf{Finding Eigenvalues}

The characteristic equation is:
\begin{align}
\mydet{\vec{A} - \lambda \vec{I}} = 0
\end{align}

Substituting $\vec{A}$,
\begin{align}
\mydet{
2-\lambda & -1 & 0\\
-1 & 2-\lambda & -1\\
0 & -1 & 2-\lambda
}= 0
\end{align}

\end{frame}

\begin{frame}{Examples}

\begin{align}
\mydet{
2-\lambda & -1 & 0\\
-1 & 2-\lambda & -1\\
0 & -1 & 2-\lambda
}
\xleftrightarrow{R_2 \to R_2 + \frac{1}{2-\lambda}R_1}
\mydet{
2-\lambda & -1 & 0\\
0 & \frac{(2-\lambda)^2 - 1}{2-\lambda} & -1\\
0 & -1 & 2-\lambda
}
\xleftrightarrow{R_3 \to R_3 + \frac{(2-\lambda)}{(2-\lambda)^2 - 1}R_2}
\mydet{
2-\lambda & -1 & 0\\
0 & \frac{(2-\lambda)^2 - 1}{2-\lambda} & -1\\
0 & 0 & \frac{(2-\lambda)^3 - 4(2-\lambda)}{(2-\lambda)^2 - 1}
}
\end{align}

\begin{align}
(2-\lambda)\Big(\frac{(2-\lambda)^2 - 1}{2-\lambda}\Big)\Big(\frac{(2-\lambda)^3 - 4(2-\lambda)}{(2-\lambda)^2 - 1}\Big) = 0
\end{align}

Simplifying,
\begin{align}
(2-\lambda)\big((2-\lambda)^2 - 2\big) = 0
\end{align}

Hence, the eigenvalues are:
\begin{align}
\lambda_1 = 2,\ \lambda_2 = 2 + \sqrt{2},\ \lambda_3 = 2 - \sqrt{2}
\end{align}

All eigenvalues are real, confirming positive definiteness.\\

\textbf{Finding Eigenvectors}

\end{frame}

\begin{frame}{Examples}

\textbf{For } $\lambda_1 = 2$:
\begin{align}
(\vec{A}-2\vec{I})\vec{v} &= 0\\
\myvec{0 & -1 & 0\\ -1 & 0 & -1\\ 0 & -1 & 0}\myvec{x\\y\\z} &= \myvec{0\\0\\0}
\end{align}
\begin{align}
\myaugvec{3}{0 & -1 & 0 & 0\\ -1 & 0 & -1 & 0\\ 0 & -1 & 0 & 0}
\xleftrightarrow{R_2 \leftrightarrow R_1}
\myaugvec{3}{-1 & 0 & -1 & 0\\ 0 & -1 & 0 & 0\\ 0 & -1 & 0 & 0}\\
\xleftrightarrow{R_3 \to R_3 - R_2}
\myaugvec{3}{-1 & 0 & -1 & 0\\ 0 & -1 & 0 & 0\\ 0 & 0 & 0 & 0}
\end{align}
Hence, $x = -z,\ y = 0 \Rightarrow \vec{v}_1 = \myvec{1\\0\\-1}$.

\end{frame}

\begin{frame}{Examples}

\textbf{For } $\lambda_2 = 2+\sqrt{2}$:
\begin{align}
(\vec{A}-(2+\sqrt{2})\vec{I})\vec{v}=0\\
\myvec{-\sqrt{2} & -1 & 0\\ -1 & -\sqrt{2} & -1\\ 0 & -1 & -\sqrt{2}}\myvec{x\\y\\z} = \myvec{0\\0\\0}
\end{align}
\begin{align}
\myaugvec{3}{-\sqrt{2} & -1 & 0 & 0\\ -1 & -\sqrt{2} & -1 & 0\\ 0 & -1 & -\sqrt{2} & 0}
\xleftrightarrow{R_2 \to R_2 - \frac{1}{\sqrt{2}}R_1}
\myaugvec{3}{-\sqrt{2} & -1 & 0 & 0\\ 0 & -\frac{1}{\sqrt{2}} & -1 & 0\\ 0 & -1 & -\sqrt{2} & 0}\\
\xleftrightarrow{R_3 \to R_3 -\sqrt{2}R_2}
\myaugvec{3}{-\sqrt{2} & -1 & 0 & 0\\ 0 & -\frac{1}{\sqrt{2}} & -1 & 0\\ 0 & 0 & 0 & 0}
\end{align}

\end{frame}

\begin{frame}{Examples}

Simplifying gives proportional relation $y = -\sqrt{2}x, \ z = x$.
Hence $\vec{v}_2 = \myvec{1\\-\sqrt{2}\\1}$.

\end{frame}

\begin{frame}{Examples}

\textbf{For } $\lambda_3 = 2-\sqrt{2}$:
Similarly, solving gives $\vec{v}_3 = \myvec{1\\\sqrt{2}\\1}$.

Thus, eigenvectors are \textbf{real and mutually orthogonal}.

\pagebreak

\textbf{Example 2 : Complex Symmetric Matrix}

(\textbf{a}) Let
\begin{align}
\vec{A} = \myvec{1 & 1+i\\ 1+i & 1}
\end{align}

Here, $\vec{A} = \vec{A}^\top$ but $\vec{A} \ne \vec{A}^*$, so $\vec{A}$ is complex symmetric but not Hermitian.

\textbf{Finding Eigenvalues}

\begin{align}
\mydet{\vec{A} - \lambda \vec{I}} = 0
\end{align}

\begin{align}
\mydet{
1-\lambda & 1+i\\
1+i & 1-\lambda
} = 0
\end{align}

\end{frame}

\begin{frame}{Examples}

\begin{align}
\mydet{
1-\lambda & 1+i\\
1+i & 1-\lambda
}
\xleftrightarrow{R_2 \to R_2 - \frac{1+i}{1-\lambda} R_1}
\mydet{
1-\lambda & 1+i\\
0 & (1-\lambda) - \frac{(1+i)^2}{1-\lambda}
}
\end{align}

\begin{align}
(1-\lambda)^2 - (1+i)^2 = 0
\end{align}

\begin{align}
(1-\lambda)^2 - 2i = 0
\end{align}

\begin{align}
1-\lambda = \pm \sqrt{2i} \quad \Rightarrow \quad \lambda = 1 \mp \sqrt{2i}
\end{align}

Hence eigenvalues are:
\begin{align}
\lambda_1 = 1 - \sqrt{2i}, \quad \lambda_2 = 1 + \sqrt{2i}
\end{align}

\end{frame}

\begin{frame}{Examples}

\textbf{Finding Eigenvectors}

\textbf{For } $\lambda_1 = 1 - \sqrt{2i}$:
\begin{align}
(\vec{A}-\lambda_1 \vec{I})\vec{v}=0\\
\myvec{\sqrt{2i} & 1+i\\ 1+i & \sqrt{2i}}\myvec{x\\y}=\myvec{0\\0}
\end{align}
\begin{align}
  \myaugvec{2}{\sqrt{2i} & 1+i & 0\\ 1+i & \sqrt{2i} & 0}
\xleftrightarrow{R_2 \to R_2 - \frac{1+i}{\sqrt{2i}}R_1}
  \myaugvec{2}{\sqrt{2i} & 1+i & 0\\ 0 & 0 & 0}
\end{align}
Hence $y = -\frac{\sqrt{2i}}{1+i}x$ and $\vec{v}_1 = \myvec{1\\ -\frac{\sqrt{2i}}{1+i}}$.

\end{frame}

\begin{frame}{Examples}

\textbf{For } $\lambda_2 = 1 + \sqrt{2i}$:
\begin{align}
(\vec{A}-\lambda_2 \vec{I})\vec{v}=0\\
\myvec{-\sqrt{2i} & 1+i\\ 1+i & -\sqrt{2i}}\myvec{x\\y}=\myvec{0\\0}
\end{align}
\begin{align}
  \myaugvec{2}{-\sqrt{2i} & 1+i & 0\\ 1+i & -\sqrt{2i} & 0}
\xleftrightarrow{R_2 \to R_2 - \frac{1+i}{-\sqrt{2i}}R_1}
  \myaugvec{2}{-\sqrt{2i} & 1+i & 0\\ 0 & 0 & 0}
\end{align}
Hence $y = \frac{\sqrt{2i}}{1+i}x$ and $\vec{v}_2 = \myvec{1\\ \frac{\sqrt{2i}}{1+i}}$.

\begin{align}
\lambda_1 &= 1 - \sqrt{2i}, \quad \vec{v}_1 = \myvec{1\\ -\,\frac{\sqrt{2i}}{1+i}}\\
\lambda_2 &= 1 + \sqrt{2i}, \quad \vec{v}_2 = \myvec{1\\ \frac{\sqrt{2i}}{1+i}}
\end{align}

The Eigen vectors are \textbf{complex}

\end{frame}

\begin{frame}{Examples}


(\textbf{b}) Let
\begin{align}
\vec{A} = \myvec{1 & i & 0\\ i & 1 & i\\ 0 & i & 1}
\end{align}

Here, $\vec{A} = \vec{A}^\top$ but $\vec{A} \ne \vec{A}^*$, so $\vec{A}$ is complex symmetric but not Hermitian.

\textbf{Finding Eigenvalues}

\begin{align}
\mydet{\vec{A} - \lambda \vec{I}} = 0
\end{align}

\begin{align}
\mydet{
1-\lambda & i & 0\\
i & 1-\lambda & i\\
0 & i & 1-\lambda
} = 0
\end{align}

\end{frame}

\begin{frame}{Examples}

\begin{align}
\mydet{
1-\lambda & i & 0\\
i & 1-\lambda & i\\
0 & i & 1-\lambda
}
\xleftrightarrow{R_2 \to R_2 - \frac{i}{1-\lambda} R_1}
\mydet{
1-\lambda & i & 0\\
0 & \frac{(1-\lambda)^2 + 1}{1-\lambda} & i\\
0 & i & 1-\lambda
}
\\
\xleftrightarrow{R_3 \to R_3 - \frac{i(1-\lambda)}{(1-\lambda)^2 + 1} R_2}
\mydet{
1-\lambda & i & 0\\
0 & \frac{(1-\lambda)^2 + 1}{1-\lambda} & i\\
0 & 0 & \frac{(1-\lambda)^3 + 2(1-\lambda)}{(1-\lambda)^2 + 1}
}
\end{align}

\begin{align}
(1-\lambda)\Big(\frac{(1-\lambda)^2 + 1}{1-\lambda}\Big)\Big(\frac{(1-\lambda)^3 + 2(1-\lambda)}{(1-\lambda)^2 + 1}\Big) = 0
\end{align}

Simplifying,
\begin{align}
(1-\lambda)\big((1-\lambda)^2 + 2\big) = 0
\end{align}

Hence eigenvalues are:
\begin{align}
\lambda_1 = 1,\ 
\lambda_2 = 1 + i\sqrt{2},\ 
\lambda_3 = 1 - i\sqrt{2}
\end{align}

\end{frame}

\begin{frame}{Examples}

\textbf{Finding Eigenvectors}

\textbf{For } $\lambda_1 = 1$:
\begin{align}
(\vec{A}-\vec{I})\vec{v}=0\\
\myvec{0 & i & 0\\ i & 0 & i\\ 0 & i & 0}\myvec{x\\y\\z}=\myvec{0\\0\\0}
\end{align}
\begin{align}
  \myaugvec{3}{0 & i & 0 & 0\\ i & 0 & i & 0\\ 0 & i & 0 & 0}
\xleftrightarrow{R_2 \leftrightarrow R_1}
  \myaugvec{3}{i & 0 & i & 0\\ 0 & i & 0 & 0\\ 0 & i & 0 & 0}
\xleftrightarrow{R_3 \to R_3 - R_2}
  \myaugvec{3}{i & 0 & i & 0\\ 0 & i & 0 & 0\\ 0 & 0 & 0 & 0}
\end{align}
Hence $y=0,\ z=-x \Rightarrow \vec{v}_1 = \myvec{1\\0\\-1}$.

\end{frame}

\begin{frame}{Examples}

\textbf{For} \(\lambda_2 = 1 + i\sqrt{2}\)
\begin{align}
\vec{A} - \lambda_2 \vec{I} 
&= \myvec{
 -i\sqrt{2} & i & 0\\
 i & -i\sqrt{2} & i\\
 0 & i & -i\sqrt{2}
}
\end{align}

Forming the augmented matrix :
\begin{align}
\myaugvec{3}{
 -i\sqrt{2} & i & 0 & 0\\
 i & -i\sqrt{2} & i & 0\\
 0 & i & -i\sqrt{2} & 0
}\\
\xleftrightarrow{R_2 \to R_2 - \frac{i}{-i\sqrt{2}} R_1}
\myaugvec{3}{
 -i\sqrt{2} & i & 0 & 0\\
 0 & i\brak{\frac{1}{\sqrt{2}} - \sqrt{2}} & i & 0\\
 0 & i & -i\sqrt{2} & 0
}
\end{align}

\begin{align}
\xleftrightarrow{R_3 \to R_3 - \frac{i}{i\brak{\frac{1}{\sqrt{2}} - \sqrt{2}}} R_2}
\myaugvec{3}{
 -i\sqrt{2} & i & 0 & 0\\
 0 & i\brak{\frac{1}{\sqrt{2}} - \sqrt{2}} & i & 0\\
 0 & 0 & 0 & 0
}
\end{align}

Let \(x = z\). From the first two equations, \(y = \sqrt{2}x\).

\end{frame}

\begin{frame}{Examples}

\begin{align}
\vec{v}_2 &= \myvec{1\\[2pt]\sqrt{2}\\[2pt]1}
\end{align}

\textbf{For} \(\lambda_3 = 1 - i\sqrt{2}\)

\begin{align}
\vec{A} - \lambda_3 \vec{I} 
&=  \myvec{
 i\sqrt{2} & i & 0\\
 i & i\sqrt{2} & i\\
 0 & i & i\sqrt{2}
}
\end{align}

Forming the augmented matrix:
\begin{align}
\myaugvec{3}{
 i\sqrt{2} & i & 0 & 0\\
 i & i\sqrt{2} & i & 0\\
 0 & i & i\sqrt{2} & 0
}
\xleftrightarrow{R_2 \to R_2 - \frac{i}{i\sqrt{2}}R_1}
\myaugvec{3}{
 i\sqrt{2} & i & 0 & 0\\
 0 & i\brak{\sqrt{2} - \frac{1}{\sqrt{2}}} & i & 0\\
 0 & i & i\sqrt{2} & 0
}
\end{align}

\end{frame}

\begin{frame}{Examples}

\begin{align}
\xleftrightarrow{R_3 \to R_3 - \frac{i}{i\brak{\sqrt{2} - \frac{1}{\sqrt{2}}}} R_2}
\myaugvec{3}{
 i\sqrt{2} & i & 0 & 0\\
 0 & i\brak{\sqrt{2} - \frac{1}{\sqrt{2}}} & i & 0\\
 0 & 0 & 0 & 0
}
\end{align}

Hence, \(y = -\sqrt{2}x\), \(x = z\).

\begin{align}
\vec{v}_3 &= \myvec{1\\[2pt]-\sqrt{2}\\[2pt]1}
\end{align}

\end{frame}

\begin{frame}{Examples}

Thus, by Gaussian elimination:
\begin{align}
\lambda_1 &= 1, & \vec{v}_1 &= \myvec{1\\0\\-1}\\
\lambda_2 &= 1 + i\sqrt{2}, & \vec{v}_2 &= \myvec{1\\\sqrt{2}\\1}\\
\lambda_3 &= 1 - i\sqrt{2}, & \vec{v}_3 &= \myvec{1\\-\sqrt{2}\\1}
\end{align}

Some of the eigen values are complex but the eigen vectors are \textbf{real}.

\end{frame}

\begin{frame}{Examples}

\textbf{(c)}

\begin{align}
\vec{A} &= \myvec{1 & i \\ 0 & 2} \\[2mm]
\mydet{\vec{A} - \lambda \vec{I}} 
&= \mydet{1-\lambda & i \\ 0 & 2-\lambda}
= (1-\lambda)(2-\lambda) \\[2mm]
\lambda_1 &= 1, \quad \lambda_2 = 2
\end{align}

For $\lambda_1 = 1$,
\begin{align}
\vec{A} - \lambda_1 \vec{I} &= \myvec{0 & i \\ 0 & 1} \\
i y &= 0 \Rightarrow y = 0 \\[1mm]
\vec{v_1} &= \myvec{1 \\ 0}
\end{align}

\end{frame}

\begin{frame}{Examples}

For $\lambda_2 = 2$,
\begin{align}
\vec{A} - \lambda_2 \vec{I} &= \myvec{-1 & i \\ 0 & 0}
\xleftrightarrow{R_1/(-1)}
\myvec{1 & -i \\ 0 & 0} \\[1mm]
x &= i y \\[1mm]
\vec{v_2} &= \myvec{i \\ 1}
\end{align}

Check orthogonality:
\begin{align}
\vec{v_1}^*\vec{v_2} \ne 0
\end{align}

Hence, the eigenvectors are not orthogonal.

\end{frame}

\end{document}
